\section{Introduction}

Resting-state functional connectomes summarize statistical associations among
regional blood-oxygen-level-dependent time series. Representing brain regions
as nodes and functional associations as edges makes graph learning an
intuitive modeling choice. Yet the common node feature in connectome studies
is itself an entire row of the connectivity matrix, so each node begins with a
global connectivity profile rather than purely local information. In this
setting, neighborhood aggregation may remove useful distinctions instead of
adding missing context.

Recent multi-dataset benchmarking reported that classical models and
non-graph neural networks can match or exceed graph deep-learning models for
functional-connectome prediction, and isolated message aggregation as a
potential source of degradation \citep{han2026rethinking}. That result changes
the relevant question. The value of a graph operator should not be presumed
from the fact that the input can be drawn as a graph; it should be demonstrated
against identity propagation and strong regularized baselines under the same
data splits.

BuNNs diffuse features after transporting them through learned orthogonal maps
between node-specific frames \citep{bamberger2025bundle}. Their theoretical and
synthetic motivation includes greater control over representation collapse
than ordinary graph diffusion. Related neural-sheaf work likewise shows that
non-trivial transport can change diffusion behavior \citep{bodnar2022sheaf}.
Whether that advantage transfers to dense functional connectomes with global
connectivity-profile features is an empirical question.

The initial formulation of this project connected negative functional
correlations with heterophilic message passing. Formalizing the experiment
showed that this interpretation was not identifiable: correlation sign does
not establish node-label heterophily, excitation, inhibition, anatomical
connectivity, or causal information flow, and sign can depend on preprocessing
choices such as global-signal regression \citep{murphy2017global}. The dataset,
BuNN architecture, connectome representation, and aggregation problem were
therefore retained, while the hypothesis was narrowed to a controlled operator
comparison.

This study asks whether learned bundle transport produces a more favorable
held-out-site performance-density curve than GCN aggregation, whether any
advantage exceeds identity and learned-local controls as well as a regularized
connectome baseline, and whether predictive behavior co-occurs with
common-frame representation preservation. The contribution is an audit of an
architectural inductive bias under one pre-specified ABIDE-I pipeline, not a
claim about biological geometry or clinical diagnosis.
